\subsection{Quasi-linear Mailbox Types}

Pat uses quasi-linear typing \cite{kobayashiQuasilinearTypes1999,ennalsLinearTypesPacket2004a} to solve the problem of name aliasing and cyclic dependency between mailboxes.
To this end, every mailbox type is equip\-ped with a \emph{usage annotation}: either \emph{second class} ($\usecondclass$) or \emph{returnable} ($\ureturnable$).
Values of second class types can be used several times within the scope they are
defined, however, they cannot leave that scope.  On the other hand, values of
returnable types can be let-bound and used as the subject of a guard, but must
appear as the last lexical occurrence of the name.
Base types are not equipped with usage, as they do not suffer from the aliasing problems.

\begin{figure}[t]
    {\small
  \centering
  \begin{tabular}{p{7.5cm}p{5cm}}
    \begin{tabular}{lccl}
      Usage annotations (\rocqlink{https://edgardschi.github.io/mailbox-types-rocq/MailboxTypes.Types.html\#UsageAnnotation}) \; & $\eta$ \; & $::=$ \; & $\usecondclass \mid \ureturnable$\\
      Usage-annotated types (\rocqlink{https://edgardschi.github.io/mailbox-types-rocq/MailboxTypes.Types.html\#TUsage}) \; & $A, B$ \; & $::=$ \; & $C \mid J^\eta$\\
    \end{tabular}
    &
    \begin{tabular}{p{2.5cm}p{2.5cm}}
      \multicolumn{2}{l}{Type operations (\rocqlink{https://edgardschi.github.io/mailbox-types-rocq/MailboxTypes.Types.html\#returnType},
        \rocqlink{https://edgardschi.github.io/mailbox-types-rocq/MailboxTypes.Types.html\#returnUsage},
        \rocqlink{https://edgardschi.github.io/mailbox-types-rocq/MailboxTypes.Types.html\#secondType},
        \rocqlink{https://edgardschi.github.io/mailbox-types-rocq/MailboxTypes.Types.html\#secondUsage})}\\
      \begin{tabular}{lcl}
        $\secondclassop{C}$ & $=$ & $C$\\
        $\secondclassop{J}$ & $=$ & $J^\usecondclass$\\
        $\secondclassop{J^\eta}$ & $=$ & $J^\usecondclass$
      \end{tabular}
      &
      \begin{tabular}{lcl}
        $\returnableop{C}$ & $=$ & $C$\\
        $\returnableop{J}$ & $=$ & $J^\ureturnable$\\
        $\returnableop{J^\eta}$ & $=$ & $J^\ureturnable$
      \end{tabular}
    \end{tabular}
  \end{tabular}

  \vspace{.5cm}

  \begin{minipage}[t]{0.53\textwidth}
    \begin{tabular*}{\linewidth}{@{}l@{\extracolsep{\fill}}r@{}}
      \textbf{Usage-annotated} & \fbox{$\mbtucomb{A}{B}{A'}$}\\
      \textbf{type combinations} (\rocqlink{https://edgardschi.github.io/mailbox-types-rocq/MailboxTypes.Types.html\#TypeUsageCombination}) &
    \end{tabular*}
    \begin{center}
      \begin{prooftree}
        \infer0{\mbtucomb{C}{C}{C}}
      \end{prooftree}
      \quad
      \begin{prooftree}
        \hypo{\mbtucomb{\eta_1}{\eta_2}{\eta}}
        \hypo{\mbtcomb{J_1}{J_2}{K}}
        \infer2{\mbtucomb{J_1^{\eta_1}}{J_2^{\eta_2}}{K^{\eta}}}
      \end{prooftree}
    \end{center}
  \end{minipage}
  \hfill
  \begin{minipage}[t]{0.46\textwidth}\vspace{-\baselineskip}
    \begin{tabular*}{\linewidth}{@{}l@{\extracolsep{\fill}}r@{}}
      \textbf{Usage} & \fbox{$\mbtucomb{\eta_1}{\eta_2}{\eta_3}$}\\
      \textbf{combinations} (\rocqlink{https://edgardschi.github.io/mailbox-types-rocq/MailboxTypes.Types.html\#UsageCombination})&
    \end{tabular*}
    \begin{mathpar}
      \mbtucomb{\usecondclass}{\usecondclass}{\usecondclass}

      \mbtucomb{\usecondclass}{\ureturnable}{\ureturnable}
  \end{mathpar}
  \end{minipage}

  \vspace{.5cm}

  \begin{minipage}[t]{0.37\textwidth}\vspace{0cm}
    \begin{tabular*}{\linewidth}{@{}l@{\extracolsep{\fill}}r@{}}
      \textbf{Usage} & \fbox{$\eta_1 \subtype \eta_2$}\\
      \textbf{subtyping} (\rocqlink{https://edgardschi.github.io/mailbox-types-rocq/MailboxTypes.Types.html\#UsageSubtype}) &
    \end{tabular*}
    \begin{mathpar}
      \eta \subtype \eta

      \ureturnable \subtype \usecondclass
    \end{mathpar}
  \end{minipage}\hfill
  \begin{minipage}[t]{0.6\textwidth}\vspace{0pt}
    \textbf{Subtyping} (\rocqlink{https://edgardschi.github.io/mailbox-types-rocq/MailboxTypes.Types.html\#Subtype}) \hfill \fbox{$A \subtype B$}
    \begin{center}
          \begin{prooftree}
              \infer0{C \subtype C}
          \end{prooftree}
          \quad
          \begin{prooftree}
              \hypo{E \mbpinc F}
              \hypo{\eta_1 \subtype \eta_2}
              \infer2{\mbtin{E^{\eta_1}} \subtype \mbtin{F^{\eta_2}}}
          \end{prooftree}
          \quad
          \begin{prooftree}
              \hypo{F \mbpinc E}
              \hypo{\eta_1 \subtype \eta_2}
              \infer2{\mbtout{E^{\eta_1}} \subtype \mbtout{F^{\eta_2}}}
          \end{prooftree}
      \end{center}
  \end{minipage}
}

  \caption{Definitions of usage-annotated types.}
  \label{fig:usage-annotation-definition}
\end{figure}

Figure \ref{fig:usage-annotation-definition} shows definitions of usage annotations, usage-annotated types and relations on them.
We also define operators $\secondclassop{\cdot}$ and $\returnableop{\cdot}$ to turn types into second class and returnable respectively.
%
Additionally, we define a mailbox type to be \emph{unrestricted} (written $\unr{A}$) (\rocqlink{https://edgardschi.github.io/mailbox-types-rocq/MailboxTypes.Types.html\#Unrestricted})
if it is a base type, or $A = \mbtout{\mbone^\usecondclass}$.

% \begin{figure}[t]
%   \begin{minipage}[t]{0.37\textwidth}\vspace{0cm}
%     % \textbf{Usage subtyping} \hfill \fbox{$\eta_1 \subtype \eta_2$}
%     \begin{align*}
%       \eta \subtype \eta\\
%       \ureturnable \subtype \usecondclass
%     \end{align*}
%   \end{minipage}\hfill
%   \begin{minipage}[t]{0.6\textwidth}\vspace{0pt}
%     % \textbf{Subtyping} \hfill \fbox{$A \subtype B$}
%     \begin{center}
%           \begin{prooftree}
%               \infer0{C \subtype C}
%           \end{prooftree}
%           \quad
%           \begin{prooftree}
%               \hypo{E \mbpinc F}
%               \hypo{\eta_1 \subtype \eta_2}
%               \infer2{\mbtin{E^{\eta_1}} \subtype \mbtin{F^{\eta_2}}}
%           \end{prooftree}
%           \quad
%           \begin{prooftree}
%               \hypo{F \mbpinc E}
%               \hypo{\eta_1 \subtype \eta_2}
%               \infer2{\mbtout{E^{\eta_1}} \subtype \mbtout{F^{\eta_2}}}
%           \end{prooftree}
%       \end{center}
%   \end{minipage}
%
%   \caption{Subtyping with usage annotations.}
%   \label{fig:usage-subtyping}
% \end{figure}

% \subsubsection{Usage-annotated Type Combination.}
\paragraph{Usage-annotated Type Combination.}
We extend type combination to \emph{usage-annotated type combination} by incorporating usage combination.
% Figure \ref{fig:usage-annotated-type-combination} shows the extension of type combination to \emph{usage-annotated type combination}.
For usage annotations, we have to obey the rule that a $\usecondclass$ variable has to occur before a $\ureturnable$ variable.
Otherwise, we would violate the property that any $\ureturnable$ variable appears as the final one in a process.
Combining two $\ureturnable$ types is not allowed as there could otherwise be multiple instances of a returnable mailbox name per thread.

\paragraph{Subtyping.}
Subtyping is needed in mailbox typing to rewrite mailbox types so that they can
be combined by the type combination operator.
For mailbox types, subtyping is \emph{covariant} for types with receive
capability and \emph{contravariant} for types with send capability
\cite{deliguoroMailboxTypesUnordered2018}.
Subtyping also allows us to treat returnable types as second-class types.

%%   Subtyping on usage annotations only
%%   allows reflexivity, or for returnable types to be subtypes of second class
%%   types.  If we would also allow subtyping in the other direction, i.e.\
%%   $\usecondclass \subtype \ureturnable$, it would be possible to escape the
%%   restrictions of second class types using subtyping.

%% SJF: I don't think this really adds much
%%   \paragraph{Mechanization.}
%%
%%   The addition of usage annotations to mailbox types makes the mechanization more
%%   complex in that proofs require reasoning about usage.  Fortunately, we can use
%%   already proven properties about non-usage counterpart definitions from previous
%%   sections.  With them, we only have to argue about changes introduced by having a
%%   restriction on usage annotations of types.

